\documentclass[11pt]{article}
\pagestyle{plain}

% \pgfplotsset{compat=1.18}
\usepackage[margin=1in]{geometry}
\usepackage{amsmath,amsfonts,amssymb,amsthm}
\usepackage[ruled,vlined,linesnumbered,commentsnumbered]{algorithm2e}
\SetKwInput{KwIn}{Input}
\SetKwInput{KwOut}{Output}
\SetKwInput{KwGlobal}{Global variable}
\SetKwInput{KwOracle}{Oracle}
\SetKwComment{Comment}{// }{}
\usepackage{graphicx}
\usepackage{enumerate}
\usepackage{bbm}
\usepackage{verbatim}
\usepackage{hyperref,color}
\usepackage[capitalize,nameinlink]{cleveref}
\usepackage[dvipsnames]{xcolor}
\hypersetup{
	colorlinks=true,
	pdfpagemode=UseNone,
	citecolor=OliveGreen,
	linkcolor=NavyBlue,
	urlcolor=Magenta,
	pdfstartview=FitW
}
\usepackage{appendix}
\usepackage{makecell}
\usepackage{tablefootnote}
\crefname{appsec}{Appendix}{Appendices}
\usepackage{tikz}
\usepackage{pgfplots}
\usepackage{xifthen}
\usepackage{subcaption}
\usepackage{hhline}
\usepackage{listings}

\theoremstyle{plain}
\newtheorem{theorem}{Theorem}[section]
\newtheorem{proposition}[theorem]{Proposition}
\newtheorem{lemma}[theorem]{Lemma}
\newtheorem{corollary}[theorem]{Corollary}
\newtheorem{conjecture}[theorem]{Conjecture}
\newtheorem{problem}[theorem]{Problem}
\newtheorem{claim}[theorem]{Claim}
\newtheorem{fact}[theorem]{Fact}

\theoremstyle{definition}
\newtheorem{definition}[theorem]{Definition}
\newtheorem{example}[theorem]{Example}
\newtheorem{observation}[theorem]{Observation}
\newtheorem{assumption}[theorem]{Assumption}
\newtheorem*{assumption*}{Assumption}
\newtheorem{condition}[theorem]{Condition}

\theoremstyle{remark}
\newtheorem{remark}[theorem]{Remark}

\crefname{lemma}{Lemma}{Lemmas}
\crefname{theorem}{Theorem}{Theorems}
\crefname{definition}{Definition}{Definitions}
\crefname{fact}{Fact}{Facts}
\crefname{claim}{Claim}{Claims}
\crefname{proposition}{Proposition}{Propositions}


\newcommand{\dif}{\,\mathrm{d}}
%\newcommand{\E}{\mathbb{E}}
%\newcommand{\Var}{\mathrm{Var}}
\newcommand{\Cov}{\mathrm{Cov}}
\newcommand{\MI}{\mathrm{I}}
%\newcommand{\Ent}{\mathrm{Ent}}
\newcommand{\ind}{\mathbbm{1}}
\newcommand{\allone}{\mathbf{1}}
\newcommand{\tr}{\mathrm{tr}}
\newcommand{\wt}{\mathrm{weight}}
\DeclareMathOperator*{\argmax}{arg\,max}
\newcommand{\norm}[1]{\left\lVert #1 \right\rVert}
\newcommand{\TV}[2]{\left\|#1 - #2\right\|_{\mathrm{TV}}}
\newcommand{\ceil}[1]{\left\lceil #1 \right\rceil}
\newcommand{\floor}[1]{\left\lfloor #1 \right\rfloor}
\newcommand{\ip}[2]{\left\langle #1 , #2 \right\rangle}
\newcommand{\grad}{\nabla}
\newcommand{\hessian}{\nabla^2}
\newcommand{\trans}{\intercal}
\newcommand{\diag}{\mathrm{diag}}
\newcommand{\poly}{\mathrm{poly}}
\newcommand{\dist}{\mathrm{dist}}
\newcommand{\sgn}{\mathrm{sgn}}
\newcommand{\fpras}{\mathsf{FPRAS}}
\newcommand{\fpaus}{\mathsf{FPAUS}}
\newcommand{\fptas}{\mathsf{FPTAS}}

\newcommand{\eps}{\varepsilon}

\newcommand{\N}{\mathbb{N}}
\newcommand{\Z}{\mathbb{Z}}
\newcommand{\Q}{\mathbb{Q}}
\newcommand{\R}{\mathbb{R}}
\newcommand{\F}{\mathbb{F}}
\newcommand{\C}{\mathbb{C}}

\renewcommand{\AA}{\mathcal{A}}
\newcommand{\BB}{\mathcal{B}}
\newcommand{\CC}{\mathcal{C}}
\newcommand{\DD}{\mathcal{D}}
\newcommand{\EE}{\mathcal{E}}
\newcommand{\FF}{\mathcal{F}}
\newcommand{\GG}{\mathcal{G}}
\newcommand{\HH}{\mathcal{H}}
\newcommand{\II}{\mathcal{I}}
\newcommand{\JJ}{\mathcal{J}}
\newcommand{\KK}{\mathcal{K}}
\newcommand{\LL}{\mathcal{L}}
\newcommand{\MM}{\mathcal{M}}
\newcommand{\NN}{\mathcal{N}}
\newcommand{\OO}{\mathcal{O}}
\newcommand{\PP}{\mathcal{P}}
\newcommand{\QQ}{\mathcal{Q}}
\newcommand{\RR}{\mathcal{R}}
\renewcommand{\SS}{\mathcal{S}}
\newcommand{\TT}{\mathcal{T}}
\newcommand{\UU}{\mathcal{U}}
\newcommand{\VV}{\mathcal{V}}
\newcommand{\WW}{\mathcal{W}}
\newcommand{\XX}{\mathcal{X}}
\newcommand{\YY}{\mathcal{Y}}
\newcommand{\ZZ}{\mathcal{Z}}

\newcommand{\KL}[2]{D_{\mathrm{KL}} \left( #1 \,\Vert\, #2 \right)}
\newcommand{\KLbig}[2]{D_{\mathrm{KL}} \left( #1 \,\big\Vert\, #2 \right)}
\newcommand{\KLBig}[2]{D_{\mathrm{KL}} \left( #1 \,\Big\Vert\, #2 \right)}

\newcommand{\SKL}[2]{D_{\mathrm{SKL}} \left( #1 \,\Vert\, #2 \right)}
\newcommand{\SKLbig}[2]{D_{\mathrm{SKL}} \left( #1 \,\big\Vert\, #2 \right)}
\newcommand{\SKLBig}[2]{D_{\mathrm{SKL}} \left( #1 \,\Big\Vert\, #2 \right)}

\newcommand{\qandq}{\quad\text{and}\quad}
\newcommand{\supp}{\mathrm{supp}}

\newcommand{\DKL}[2]{\-D_{\-{KL}}\left(#1 \parallel #2\right)}
\newcommand{\DTV}[2]{\-D_{\mathrm{TV}}\left({#1},{#2}\right)}
% \newcommand{\dist}{\mathrm{dist}}
\newcommand{\e}{\mathrm{e}}
\renewcommand{\epsilon}{\varepsilon}
\renewcommand{\emptyset}{\varnothing}
% \newcommand{\norm}[1]{\left\Vert#1\right\Vert}
\newcommand{\set}[1]{\left\{#1\right\}}
\newcommand{\tuple}[1]{\left(#1\right)} 
% \newcommand{\eps}{\varepsilon}
\newcommand{\inner}[2]{\left\langle #1,#2\right\rangle}
\newcommand{\tp}{\tuple}
\newcommand{\ol}{\overline}
\newcommand{\abs}[1]{\left\vert#1\right\vert}
\newcommand{\ctp}[1]{\left\lceil#1\right\rceil}
\newcommand{\ftp}[1]{\left\lfloor#1\right\rfloor}
\newcommand{\Ind}[1]{\-{Ind}_{#1}}
\newcommand{\todo}[1]{{\color{red} TODO: #1}}

\newcommand{\Cor}{\Psi^{\mathrm{sym}}}
\newcommand{\cor}{\Psi^{\-{cor}}}
\newcommand{\Inf}{\Psi}
% \newcommand{\diag}{\-{diag}}
\newcommand{\Tr}{\-{Tr}}
\newcommand{\rank}{\-{rank}}

\def\*#1{\boldsymbol{#1}} % Use \*A for \mathbf{A}
\def\+#1{\mathcal{#1}} % Use \+A for \mathcal{A}
\def\-#1{\mathrm{#1}} % Use \-A for \mathrm{A}
\def\=#1{\mathbb{#1}} % Use \^A for \mathbb{A}
\def\!#1{\mathfrak{#1}} % Use \!A for \mathfrak{A}

\newcommand*\diff{\mathop{}\!\mathrm{d}}

% probability
% \DeclareMathOperator*{\oPr}{\mathbf{Pr}}
\def\oPr{\mathop{\mathrm{Pr}}}
\renewcommand{\Pr}[2][]{ \ifthenelse{\isempty{#1}}
  {\oPr\left[#2\right]}
  {\oPr_{#1}\left[#2\right]} } % Use \Pr[a]{b} for \mathbf{Pr}_a[b], \Pr{b} for  \mathbf{Pr}[b]

% expectation: relies on the package "dsfont"
% \DeclareMathOperator*{\oE}{\mathbb{E}}
\def\oE{\mathop{\mathbb{E}}}
\newcommand{\E}[2][]{ \ifthenelse{\isempty{#1}}
  {\oE\left[#2\right]}
  {\oE_{#1}\left[#2\right]} }

% variance
%\DeclareMathOperator*{\oVar}{\mathbf{Var}}
\def\oVar{\mathrm{Var}}
\newcommand{\Var}[2][]{ \ifthenelse{\isempty{#1}}
  {\oVar\left[#2\right]}
  {\oVar_{#1}\left[#2\right]} }

% entropy
% \DeclareMathOperator*{\oEnt}{\mathbf{Ent}}
\def\oEnt{\mathrm{Ent}}
\newcommand{\Ent}[2][]{ \ifthenelse{\isempty{#1}}
  {\oEnt\left[#2\right]}
  {\oEnt_{#1}\left[#2\right]} }

\newcommand{\PhiEnt}[2][]{ \ifthenelse{\isempty{#1}}
  {\oEnt^\phi\left[#2\right]}
  {\oEnt^\phi_{#1}\left[#2\right]} }


\newenvironment{Exampleparagraph}[1][]
{
    \par 
    \noindent 
    \textbf{\color{blue} #1}
    \par 
    \vspace{0.5em} 
    \itshape 
}
{
    \par 
    \vspace{1em} 
}

%\input{local-macro}

\newcommand{\RED}[1]{\textcolor{red}{#1}}

\newcommand{\bern}[1]{\mathrm{Bern}\left(#1\right)}
\newcommand{\pll}{\parallel}



\title{MATH 6112. FALL 2026. ASSIGNMENT 3.}
\date{Due on November 6, 2026}
\author{{\textbf{Zejia Chen}} \\ \textbf{GTID: 904178159}}

\pgfplotsset{compat=1.18} 
\begin{document}
\maketitle
\section*{Problem 1}
\textbf{(a)} We have that the product of two Hermitian matrices may not be Hermitian.
$$
    \begin{pmatrix}
        1 & 1
        \\1 & 1
    \end{pmatrix}\begin{pmatrix}
        1 & 0
        \\0 & 2
    \end{pmatrix}
    = \begin{pmatrix}
        1 & 2
        \\1 & 2
    \end{pmatrix} \neq \begin{pmatrix}
        1 & 1
        \\2 & 2
    \end{pmatrix}
$$
\textbf{(b)}
Define 
$$
    A = \begin{pmatrix}
        2 & 0
        \\0 & 10
    \end{pmatrix}, 
    B = \begin{pmatrix}
        1 & 2
        \\2 & 5
    \end{pmatrix}.
$$
Trivially $A\succ 0$.
We have that $B_{11} = \det(B) = 1 > 0$ and thus, $B\succ 0$.
Since $(A-B)_{11} = 1 > 0$ and $\det(A-B) = 1 > 0$, $A\succ B$.
However, $(A^2 - B^2)_{11} = -1 < 0$ and hence, $A^2 \nsucc B^2$.
This is a valid counterexample.
\\\textbf{(c)}
Define 
$$
    A = \begin{pmatrix}
        1 & 1
        \\1 & 2
    \end{pmatrix}, 
    B = \begin{pmatrix}
        1 & -2
        \\-2 & 5
    \end{pmatrix}.
$$
We have that $\det(B) = \det(A) = A_{11} = B_{11} = 1 > 0$ and thus $A\succ 0$, $B\succ 0$.
However, $(AB + BA)_{11} = -2 < 0$ and hence, $AB+BA \nsucc 0$.
This is a valid counterexample.
\\\textbf{(d)}
Define $B = (1\ 0)^\top$ and $C = (1\ 2)^\top$.
We have that
$$
    CB^\top = \begin{pmatrix}
    1 & 0
    \\2 & 0
    \end{pmatrix} =: H,\quad H^\top H = \begin{pmatrix}
    5 & 0
    \\0 & 0
    \end{pmatrix}
    \implies \Sigma(CB^\top) = \left\{\sqrt 5, 0\right\},
    \qquad 
    BC^\top = (1) .
$$
This gives a valid counterexample.
\\\textbf{(e)}Define
$$
    A = \begin{pmatrix}
        2 & 0
        \\0 & 1
    \end{pmatrix},
    B =\begin{pmatrix}
        1 & 0
        \\0 & 2
    \end{pmatrix}
    \implies A - B = \begin{pmatrix}
        1 & 0
        \\0 & -1
    \end{pmatrix} .
$$
Then $A-B$ is clearly not positive semidefinite.
\section*{Problem 2}
Since $A,B$ are Hermitian, there exists unitary matrices $U,V$ such that 
$$
    A = U^* \Lambda_A U\qquad\text{and}\qquad B = V^*\Lambda_B V,
$$
where $\Lambda_A = \diag\{\lambda_i(A)\}_{i\in[n]}$ and $\Lambda_B = \diag\{\lambda_i(B)\}_{i\in[n]}$.
Therefore, we have that
$$
    V^*UAU^*V - B = V^*(UAU^*-\Lambda_B)V = V^*(\Lambda_A-\Lambda_B)V.
$$
Let $W = U^*V$ and then $W$ is also unitary.
Since $\lambda_i(A)\geq \lambda_i(B)$ for all $i\in[n]$, $\Lambda_A-\Lambda_B\succeq 0$ and hence, $ W^*AW \succeq B$.
\section*{Problem 3}
By the singular value decomposition, there exists unitary matrices $U,V$ such that $A= U\Sigma V$, where $\Sigma = \diag\{\sigma_i\}_{i\in[n]}$.
Let $W = (w_{ij})_{i,j\in[n]} = VU$.
Then $W$ is unitary.
Therefore, we have that
$$
    \abs{\tr(A)} = \abs{\tr(U\Sigma V)} = \abs{\tr(W\Sigma)}
    \leq \sum_{i=1}^n \sigma_i\abs{w_{ii}} \leq \sum_{i=1}^n \sigma_i,
$$
where the last inequality follows from $\abs{w_{ii}}\leq 1$ since $W$ is unitary.
\section*{Problem 4}
Let $n$ be the size of the matrix.
Observe that 
$$
    M = \begin{pmatrix}
        A^{1/2} & 0
        \\0 & B^{1/2}
    \end{pmatrix}\begin{pmatrix}
        I_n & I_n
        \\I_n & I_n
    \end{pmatrix}\begin{pmatrix}
        A^{1/2} & 0
        \\0 & B^{1/2}
    \end{pmatrix} \implies \-{rank}(M) \leq \-{rank}\begin{pmatrix}
        I_n & I_n
        \\I_n & I_n
    \end{pmatrix} = n .
$$
Therefore, $\-{rank}(M) < 2n$ and $M$ is not full rank, implying $\det(M) = 0$.
Also, for any $x_1,x_2\in \C^n$, we have that
$$
    \begin{pmatrix}
        x_1^* & x_2^*
    \end{pmatrix}\begin{pmatrix}
        A & A^{1/2}B^{1/2}
        \\B^{1/2}A^{1/2} & B
    \end{pmatrix}\begin{pmatrix}
        x_1 \\ x_2
    \end{pmatrix}
    = \norm{A^{1/2}x_1 + B^{1/2}x_2}_{2}^2 \geq 0.
$$
Hence, $M\succeq 0$.

WLOG, assume $A$ is positive definite.
Then $A^{1/2}$ is also positive definite.
Direct calculation gives
$$
    \forall x_2\in \C^n, \quad\begin{pmatrix}
        A & A^{1/2}B^{1/2}
        \\B^{1/2}A^{1/2} & B
    \end{pmatrix}\begin{pmatrix}
        -A^{-1/2}B^{1/2}x_2
        \\x_2
    \end{pmatrix}  = 0
$$
$$
    \begin{pmatrix}
        A & A^{1/2}B^{1/2}
        \\B^{1/2}A^{1/2} & B
    \end{pmatrix} \begin{pmatrix}
        x_1
        \\x_2
    \end{pmatrix} = 0
    \implies 
    A^{1/2}(A^{1/2}x_1 + B^{1/2}x_2) = 0
    \implies x_1 = -A^{-1/2}B^{1/2}x_2 .
$$
Therefore, $\ker(M) = \{(x_1^*\ x_2^*)^*\mid x_1 = -A^{-1/2}B^{1/2}x_2\}$ and the null space of $M$ has dimension $n$.
\section*{Problem 5}
Let $M(A,B)$ be the matrix defined in Problem \textbf{(4)}.
We have that $M(A,B)\succeq 0$ (``at least one of them be positive definite'' is condition is redundant for $M(A,B)\succeq 0$).
Define $M'(\eps) = M(A,B) + M(B,A) + \eps I_{2n}$ and thus, $M' \succ 0$ for $\eps > 0$.
We write $C = A+B+\eps I\succ 0$ and $D = A^{1/2}B^{1/2}+B^{1/2}A^{1/2}$.
Define $S = \begin{pmatrix}
    I & X
    \\0 & I
\end{pmatrix}$, where $X = - C^{-1}D$.
Therefore, we have that by definition of $X$,
$$
    S^* M'(\eps) S = \begin{pmatrix}
    C & CX+D
    \\X^*C+D & C+X^*CX+X^*D+DX
\end{pmatrix}=\begin{pmatrix}
    C & 0
    \\0 & C-DC^{-1}D
\end{pmatrix}.
$$
Since the congruent transformation preserves the positive definiteness, we have that $S^* M'(\eps) S \succ 0$ and thus, $C-DC^{-1}D\succ 0$.
By congruent transformation, 
$$
    C-DC^{-1}D\succ 0 \implies I\succ \tp{C^{-1/2}DC^{-1/2}}^2 .
$$
Let $H = C^{-1/2}DC^{-1/2}$ and then $H = H^*$.
Therefore, $H^2 = HH^* \succeq 0$ and we have that
$$
    I \succ H^2 \succeq 0 \implies 1 \geq \det(H)^2
    \iff \det(C)^2 \geq \det(D)^2.
$$
Since $C\succ 0$ and hence $\det(C)>0$, we have that $\det(C)\geq \det(D)$.
Letting $\eps$ tend to $0$, we conclude that
$$
    \det(A+B) \geq \det\tp{A^{1/2}B^{1/2}+B^{1/2}A^{1/2}} .
$$
\section*{Problem 6}
For any $x\in \R^n$, we have that 
$$
    x^\top (A\circ B)x = \sum_{i=1}^n \sum_{j=1}^n x_i a_{ij}b_{ij}x_j
    = \Tr\tp{D_x A D_x B^\top}
    = \Tr\tp{B^{1/2}D_x A^{1/2}A^{1/2}D_xB^{1/2}} .
$$
Let $ M = (m_{ij})_{i,j\in[n]} = B^{1/2}D_x A^{1/2}$. Then $M^\top = A^{1/2}D_xB^{1/2}$.
Therefore,
$$
    x^\top (A\circ B)x = \Tr\tp{MM^\top} = \sum_{i=1}^n \sum_{j=1}^n m_{ij}^2 \geq 0,
$$
which implies that $A\circ B$ is positive semidefinite.

Furthermore, if $A\succ 0$ and $B\succ 0$, then both $A^{1/2}$ and $B^{1/2}$ are full rank.
If $x\neq \*0$, then $\-{rank}(D_x) \geq 1$ and therefore,
$$
    \rank(M) = \rank(B^{1/2}D_xA^{1/2}) = \rank(D_x) \geq 1,
$$
implying that $M\neq 0$ and thus, $\Tr\tp{MM^\top}\neq 0$.
In this case, $x^\top (A\circ B)x > 0$ and $A\circ B\succ 0$.
\section*{Acknowledgement}
I used ChatGPT to help construct and check the counterexamples in Problems \textbf{1(b)} and \textbf{1(c)}.
\end{document}